\documentclass[11pt]{article}
% Design of the jehlp.net minimal papers, copied from the Erdos 1016 papers
% (erdos1016/paper/minimal/style/preamble.tex): a designed title block, a short list of the
% results at the start, coloured headings, boxed statements, idea boxes, and annotated figures.
%
% A paper sets \papertitle, \papersubtitle and \paperrunning before \begin{document}
% and calls \maketitleblock after it.

\usepackage[letterpaper,left=1.12in,right=1.12in,top=1.05in,bottom=1.1in]{geometry}
\usepackage{amsmath,mathtools,amssymb,amsthm}
\usepackage{newpxtext}
\usepackage[varbb]{newpxmath}
\usepackage{fontspec}
\setsansfont{texgyreheros}[Extension=.otf, UprightFont=*-regular, BoldFont=*-bold,
  ItalicFont=*-italic, BoldItalicFont=*-bolditalic, Scale=0.88]
\usepackage{microtype}
\usepackage{parskip}
\usepackage{xcolor}
\usepackage{booktabs,array}
\usepackage{enumitem}
\usepackage{multicol}
\usepackage[titles]{tocloft}
\usepackage{tikz}
\usetikzlibrary{decorations.markings,arrows.meta,positioning,calc}
\usepackage[most]{tcolorbox}
\usepackage{fancyhdr}
\usepackage{titlesec}
\usepackage[font=small,labelfont={bf,sf},labelsep=period,skip=6pt]{caption}
\usepackage{url}

\definecolor{accent}{HTML}{1B4F72}
\definecolor{warm}{HTML}{C0562A}
\definecolor{slate}{HTML}{56636E}
\definecolor{ink}{HTML}{1E1E1E}

\usepackage[colorlinks=true,linkcolor=accent,citecolor=accent,urlcolor=accent]{hyperref}

% Lean identifiers: typewriter, breakable at dots and underscores.
\DeclareUrlCommand\lean{\urlstyle{tt}}

% ---------- layout ----------
\linespread{1.06}
\setlength{\parskip}{0.55em plus 0.1em minus 0.05em}
\setlist{itemsep=0.25em, topsep=0.2em, parsep=0.1em}
\titleformat{\section}{\sffamily\bfseries\Large\color{accent}}{\thesection}{0.7em}{}
\titleformat{\subsection}{\sffamily\bfseries\normalsize\color{ink}}{\thesubsection}{0.6em}{}
\titleformat{\paragraph}[runin]{\sffamily\bfseries\small\color{accent}}{}{0pt}{}[\enspace]
\titlespacing*{\section}{0pt}{22pt}{5pt}
\titlespacing*{\subsection}{0pt}{13pt}{1pt}
\titlespacing*{\paragraph}{0pt}{9pt}{0pt}
\pagestyle{fancy}
\fancyhf{}
\renewcommand{\headrulewidth}{0pt}
\fancyhead[L]{\sffamily\footnotesize\color{slate}\paperrunning}
\fancyhead[R]{\sffamily\footnotesize\color{slate}\thepage}
\setlength{\headheight}{14pt}
\setcounter{tocdepth}{1}
\setlength{\cftbeforesecskip}{2pt}
\renewcommand{\cftsecfont}{\small}
\renewcommand{\cftsecpagefont}{\small\color{slate}}
\renewcommand{\cftsecleader}{\cftdotfill{\cftdotsep}}
\setlength{\cftsecnumwidth}{1.7em}
\setlength{\emergencystretch}{2em}
\numberwithin{equation}{section}

% ---------- title block ----------
\newcommand{\papertitle}{}
\newcommand{\papersubtitle}{}
\newcommand{\paperrunning}{}
\newcommand{\paperauthors}{KNT\enspace\textperiodcentered\enspace Fable~5.2}
\newcommand{\paperdate}{}
\newcommand{\maketitleblock}{%
  \thispagestyle{plain}%
  {\sffamily\bfseries\color{accent}\fontsize{25}{29}\selectfont\raggedright\papertitle\par}%
  \vspace{5pt}%
  {\sffamily\Large\color{slate}\raggedright\papersubtitle\par}%
  \vspace{8pt}%
  {\color{accent}\hrule height 1.2pt}%
  \vspace{7pt}%
  {\sffamily\small\color{ink}\paperauthors\hfill\color{slate}\paperdate\par}%
  \vspace{10pt}}

% Small headings inside the front matter.
\newcommand{\fronthead}[1]{\par\vspace{6pt}{\sffamily\bfseries\normalsize\color{ink}#1}\par\vspace{-2pt}}
\makeatletter
\newcommand{\contentsblock}{%
  \par\noindent\begin{minipage}{\linewidth}%
  \fronthead{Contents}%
  \setlength{\parskip}{0pt}\setlength{\columnsep}{2.2em}\raggedright
  \begin{multicols}{2}
  \@starttoc{toc}
  \end{multicols}
  \end{minipage}\par}
\makeatother

% ---------- statements ----------
% Theorems, lemmas, propositions and corollaries share one counter.
% Each statement is in a tinted box with an accent rule at the left.
\newtheoremstyle{boxed}{0pt}{0pt}{\normalfont}{0pt}{\sffamily\bfseries\small\color{accent}}{.}{0.6em}{}
\theoremstyle{boxed}
\newtheorem{theorem}{Theorem}[section]
\newtheorem{lemma}[theorem]{Lemma}
\newtheorem{proposition}[theorem]{Proposition}
\newtheorem{corollary}[theorem]{Corollary}

\tcbset{guidebox/.style={enhanced, parbox=false, boxrule=0pt, frame hidden,
  arc=1.5pt, left=11pt, right=10pt, top=7pt, bottom=7pt, before skip=11pt, after skip=9pt}}
% Not breakable: a label at the start of a box that breaks at the top of a page can be lost.
\tcbset{statementbox/.style={guidebox, colback=accent!6,
  borderline west={3pt}{0pt}{accent}}}
\tcolorboxenvironment{theorem}{statementbox}
\tcolorboxenvironment{lemma}{statementbox}
\tcolorboxenvironment{proposition}{statementbox}
\tcolorboxenvironment{corollary}{statementbox}

% The main result of a paper: a framed box with a title bar.
% \begin{maintheorem}{name} uses the counter of the theorems, and \label works after it.
\newtcolorbox{mainresult}[1]{enhanced, parbox=false, colback=white, colframe=accent, boxrule=0.9pt,
  arc=1.5pt, left=11pt, right=11pt, top=8pt, bottom=8pt, title={#1},
  fonttitle=\sffamily\bfseries\small, colbacktitle=accent, coltitle=white,
  before skip=12pt, after skip=12pt}
\newenvironment{maintheorem}[1]%
  {\refstepcounter{theorem}\begin{mainresult}{Theorem~\thetheorem\ (#1)}}%
  {\end{mainresult}}

% ---------- proofs ----------
% A proof has a thin rule at the left, so that the eye finds the statements.
\renewcommand{\proofname}{\upshape\sffamily\bfseries\small\color{slate}Proof}
\renewcommand{\qedsymbol}{\textcolor{slate}{\rule{0.55em}{0.55em}}}
\tcolorboxenvironment{proof}{enhanced, breakable, parbox=false, boxrule=0pt, frame hidden,
  colback=white, borderline west={1.2pt}{0pt}{slate!35}, left=11pt, right=0pt, top=1pt, bottom=2pt,
  before skip=6pt, after skip=11pt}
% A step or a case inside a proof.
\newcommand{\step}[1]{{\sffamily\bfseries\small\color{slate}#1}}

% ---------- idea, notation and fine-print boxes ----------
\newtcolorbox{idea}[1]{guidebox, breakable, colback=warm!6, borderline west={3pt}{0pt}{warm!80},
  before upper={{\sffamily\bfseries\footnotesize\color{warm!75!black}IDEA\enspace\textbar\enspace}{\sffamily\bfseries\small #1}\par}}
\newtcolorbox{notation}[1]{guidebox, breakable, colback=black!4, borderline west={3pt}{0pt}{slate!60},
  fontupper=\small,
  before upper={{\sffamily\bfseries\footnotesize\color{slate}NOTATION\enspace\textbar\enspace}{\sffamily\bfseries\small #1}\par}}
\newtcolorbox{fineprint}{guidebox, breakable, colback=white, borderline west={1.5pt}{0pt}{black!25},
  fontupper=\footnotesize\color{black!78}, left=9pt, top=1pt, bottom=1pt,
  before upper={{\sffamily\bfseries Fine print.}\enspace}}

% ---------- figures ----------
\tikzset{
  vtx/.style={circle, fill=ink, inner sep=1.35pt},
  port/.style={circle, draw=accent, fill=white, line width=0.8pt, inner sep=1.5pt},
  edge/.style={line width=0.55pt, draw=ink},
  hot/.style={line width=1.1pt, draw=warm},
  cool/.style={line width=1.1pt, draw=accent},
  region/.style={draw=accent, fill=accent!7, line width=0.6pt, rounded corners=6pt},
  note/.style={font=\sffamily\footnotesize, text=slate},
}

% ---------- macros of the mathematics ----------
\newcommand{\EE}{\mathbb{E}}
\newcommand{\PP}{\mathbb{P}}


\renewcommand{\papertitle}{Four decisions in Cookie Clicker}
\renewcommand{\papersubtitle}{The order of purchases, the bank for golden cookies, the order of
permanent upgrades, and the time of a reset}
\renewcommand{\paperdate}{13 November 2026}
\renewcommand{\paperrunning}{Four decisions in Cookie Clicker}
\hypersetup{pdftitle={Four decisions in Cookie Clicker},pdfauthor={KNT, Fable 5.2}}

% An example has the counter of the statements and no box.
\theoremstyle{boxed}
\newtheorem{example}[theorem]{Example}

\begin{document}
\maketitleblock

\emph{Cookie Clicker} is a game in which cookies buy items that bake more cookies. Four
decisions return through the whole game: which item to buy first, how many cookies to keep for
the golden cookies, in which order to buy the permanent upgrades, and when to reset the game for
a permanent bonus. For each decision we state a model and prove the rule that is exact in that
model. Two examples give cases in which a rule loses time: the index rule with three purchases,
and the ratio rule for upgrades if the cost of a chip changes.
Each statement is about its model. The fastest way through the whole game is not known.
Fable~5.2 is a language model. Section~\ref{sec:status} states the part of each author and the
statements that Lean has checked.

\fronthead{The results in six sentences}
\begin{enumerate}[leftmargin=1.6em]
\item Let two items be too dear for the bank. The faster order has the smaller index
$E=p/c+p/\Delta$ first, where $p$ is the price, $\Delta$ is the income that the item adds, and
$c$ is the present income (Theorem~\ref{thm:pair}).
\item With three items, the smallest index first can lose time. For a finite list, a recursion
over the sets of bought items gives the least time, and a schedule with longer waits is never
faster (Theorem~\ref{thm:list}).
\item A \emph{Lucky} golden cookie pays 15\% of the bank, up to a cap. Below the cap, a
purchase increases the expected income exactly when its payback time $p/\Delta$ is less than
$\theta=1/(0.15\,g\lambda)$ (Theorem~\ref{thm:bank}).
\item If each chip costs the same effort, the fastest orders of the permanent upgrades are the
orders along which $(1-1/m)/p$ does not increase. Here $p$ is the price in chips and $m$ is the
factor of the upgrade (Theorem~\ref{thm:ratio}).
\item A pause of $D$ seconds that multiplies all later rates by $M$ shortens a task of $H$
seconds exactly when $D<H(1-1/M)$ (Proposition~\ref{prop:pause}).
\item Let a run bake cookies in proportion to $t^{\beta}$, with a speed in proportion to the
prestige. A chain of resets is fastest exactly when each run multiplies the cookies of all runs
by $v^{\star}$, the one root of an equation in $\beta$ (Theorem~\ref{thm:reset}).
\end{enumerate}

\contentsblock

\newpage
\section{The game and the models}\label{sec:intro}

\subsection{The game}

The player owns \emph{buildings} and \emph{upgrades}. They fix the number of cookies that the
game adds to the \emph{bank} in each second. We call this number the \emph{income} and write $c$
for it. A purchase takes its price from the bank. Copy number $n+1$ of a building costs
$1.15^{\,n}$ times the first copy, so the next purchase always needs a longer wait or a larger
income.

A \emph{golden cookie} appears at random times. A click on it gives one of several effects.
Section~\ref{sec:bank} is about the effect \emph{Lucky}, which pays a part of the bank.

The player can \emph{reset} the game at any time. The game calls this an ascension. A reset
removes the bank, the buildings and the upgrades. A \emph{run} is the play between two resets.
Let $A$ be the number of cookies that all runs have baked. The \emph{prestige} is
\begin{equation}\label{eq:prestige}
  P=\Bigl(\frac{A}{10^{12}}\Bigr)^{1/3}.
\end{equation}
The game uses the whole part of $P$. After a reset, each unit of prestige adds 1\% to the
income, and each new unit gives one \emph{heavenly chip}. Chips buy \emph{permanent upgrades},
which stay after a reset. All statements about the game are for version~2.053~\cite{Game}.

\subsection{The models}

Each section fixes a model: a small set of rules in which the decision has an exact answer. The
models leave out parts of the game. The game rounds prices to whole numbers. Temporary effects
change the income. Golden cookies arrive at random times.

A random income needs care, because a mean income does not give a mean time. Suppose that the
income is $1$ or $3$, each with chance one half, and that the player needs $3$ cookies. The mean
income is $2$, and at this income the wait is $\frac32$ seconds. The wait is in fact $3$ seconds
or $1$ second, and its mean is $2$ seconds. For this reason Section~\ref{sec:bank} compares
expected incomes at one moment, and it makes no statement about the time to a goal.

A guide for players and the mod Cookie Monster+ use the four rules together~\cite{Repo}. Such a
player is a heuristic. We have no proof that it is the fastest player, and no bound for its
loss.

\subsection{Authorship and verification}\label{sec:status}

KNT is a person. Fable~5.2 is the language model Claude Fable~5.2. KNT set the task and directed
the work. An earlier version of this paper, of October 2026, had the author line Claude
(Anthropic). It called the rules optimal play. An audit of that version found claims without a
proof and gave Examples~\ref{ex:three} and~\ref{ex:effort}. Fable~5.2 wrote this version. The
exact mean of Lemma~\ref{lem:interval} and the uniqueness in Theorem~\ref{thm:reset} are new in
it.

Lean~4 has checked three statements, with rational numbers. The names are in the namespace
\lean{CookieClicker}. Theorem~\ref{thm:list} is \lean{complete_global_optimal}. The times of
Example~\ref{ex:three} are \lean{greedy_counterexample}, with
\lean{counterexample_greedy_choices} for the indices. The inequality of
Proposition~\ref{prop:pause} is \lean{achievement_break_even}. The Lean kernel accepts the
proofs with the axioms \lean{propext}, \lean{Classical.choice} and \lean{Quot.sound} only. The
files are in the folder \texttt{paper/lean} of~\cite{Repo}.

The other statements have no formal proof. A program computes each number of
Tables~\ref{tab:interval} and~\ref{tab:vstar}, and the tests of~\cite{Repo} compare the numbers
with this text. We cite \cite{Charnov} and \cite{Smith} for the names of two rules, and we did
not read them for this paper. No referee has read the text.

\begin{notation}{symbols}
\begin{itemize}[leftmargin=1.2em]
\item $c$ is the income in cookies per second, and $b$ is the number of cookies in the bank.
\item An item has a price $p$ and a \emph{gain} $\Delta$: the purchase takes $p$ from the bank
and adds $\Delta$ to the income.
\item $A$ is the number of cookies that all runs have baked, and $P$ is the prestige
\eqref{eq:prestige}.
\item $\beta$ is the \emph{growth exponent} of a run: the run bakes cookies in proportion to
$t^{\beta}$. We put $\gamma=3/\beta-1$.
\item $\EE$ is an expected value.
\end{itemize}
\end{notation}

\section{The order of purchases}\label{sec:purchases}

\paragraph{Model.} The income is $c>0$ and the bank holds $b\ge0$ cookies. There is a finite
list of items. The player can buy an item when the bank holds its price. The purchase takes no
time. Between two purchases the income does not change. Each item is bought once.

A \emph{schedule} is an order $i_1,\dots,i_n$ of the items with times $0\le t_1\le\dots\le t_n$
such that the bank holds the price of $i_k$ at the time $t_k$, for each $k$. Its \emph{time} is
$t_n$. The \emph{earliest schedule} of an order buys each item at the first moment at which the
bank holds its price.

\subsection{Two purchases}

\begin{maintheorem}{the index for two purchases}\label{thm:pair}
Let $A$ and $B$ be items with positive prices and positive gains, and let
$b\le\min(p_A,p_B)$. Let $T_{AB}$ be the time of the earliest schedule with $A$ first, and
$T_{BA}$ the time with $B$ first. Put
\[
  E_i=\frac{p_i}{c}+\frac{p_i}{\Delta_i}.
\]
Then
\[
  T_{AB}-T_{BA}=\frac{\Delta_A\,\Delta_B}{(c+\Delta_A)(c+\Delta_B)}\,\bigl(E_A-E_B\bigr).
\]
So the order with $A$ first is faster exactly when $E_A<E_B$. The two orders end with the same
income and an empty bank.
\end{maintheorem}

\begin{proof}
With $A$ first, the bank reaches $p_A$ after $(p_A-b)/c$ seconds, and the purchase empties it.
The income is then $c+\Delta_A$, and the bank reaches $p_B$ after $p_B/(c+\Delta_A)$ more
seconds. The other order is symmetric. So
\[
  T_{AB}=\frac{p_A-b}{c}+\frac{p_B}{c+\Delta_A},\qquad
  T_{BA}=\frac{p_B-b}{c}+\frac{p_A}{c+\Delta_B}.
\]
The bank $b$ cancels in the difference:
\[
  T_{AB}-T_{BA}
  =p_A\Bigl(\frac1c-\frac1{c+\Delta_B}\Bigr)-p_B\Bigl(\frac1c-\frac1{c+\Delta_A}\Bigr)
  =\frac{p_A\,\Delta_B}{c\,(c+\Delta_B)}-\frac{p_B\,\Delta_A}{c\,(c+\Delta_A)}.
\]
Take the factor $\Delta_A\Delta_B/((c+\Delta_A)(c+\Delta_B))$ out of the two terms. The first
term leaves $p_A(c+\Delta_A)/(c\,\Delta_A)=E_A$, and the second leaves $E_B$.
\end{proof}

The first term of $E_i$ is the time to save for the item from an empty bank. The second term is
the \emph{payback time}: the time in which the item earns its price. The add-on Cookie
Monster~\cite{CM} shows the \emph{payback period} $\max(p_i-b,0)/c+p_i/\Delta_i$. If the bank
holds no price, this is $E_i-b/c$, and it gives the same order as $E_i$.

The condition $b\le\min(p_A,p_B)$ is necessary. Suppose that the bank holds each price and not
the sum: $\max(p_A,p_B)\le b<p_A+p_B$. Then the first purchase needs no wait, and
$T_{AB}=(p_A+p_B-b)/(c+\Delta_A)$. So the faster order has the larger gain first.

\subsection{Three purchases}

The index gives a rule for a longer list: buy the item with the smallest index, calculate the
indices again at the new income, and continue. We call this the \emph{index rule}.

\begin{example}[the index rule can lose time]\label{ex:three}
Let $b=0$ and $c=1$. Three items have these prices and gains:
\[
  A=(3,1),\qquad B=(5,4),\qquad C=(6,6).
\]
The indices are $6$, $\frac{25}{4}$ and $7$, so the index rule buys $A$. At the income $2$ the
indices of $B$ and $C$ are $\frac{15}{4}$ and $4$, so it buys $B$ and then $C$. Its time is
\[
  3+\frac52+\frac66=\frac{13}{2}.
\]
The order $B$, $C$, $A$ has the time
\[
  5+\frac65+\frac3{11}=\frac{356}{55}=\frac{13}{2}-\frac{3}{110}.
\]
The six orders have the times
$\frac{13}{2}$ ($ABC$), $\frac{53}{8}$ ($ACB$), $\frac{33}{5}$ ($BAC$), $\frac{356}{55}$
($BCA$), $\frac{395}{56}$ ($CAB$) and $\frac{538}{77}$ ($CBA$). The order $BCA$ is the fastest
(Figure~\ref{fig:three}).
\end{example}

\begin{figure}[ht]
\centering
\begin{tikzpicture}[x=1.9cm,y=1cm]
  % time axis
  \draw[edge, slate!70] (0,-2.15) -- (6.75,-2.15);
  \foreach \t in {0,...,6} {
    \draw[edge, slate!70] (\t,-2.15) -- (\t,-2.25);
    \node[note] at (\t,-2.5) {\t};
  }
  \node[note, anchor=west] at (6.8,-2.15) {seconds};
  % index rule: A, B, C
  \draw[cool] (0,0) -- (6.5,0);
  \foreach \t/\name in {3/A,5.5/B,6.5/C} {
    \node[port] at (\t,0) {};
    \node[text=accent, font=\sffamily\small\bfseries] at (\t,0.36) {\name};
  }
  \node[note] at (1.5,-0.32) {income 1};
  \node[note] at (4.25,-0.32) {income 2};
  \node[note] at (6.0,-0.32) {income 6};
  \node[note, anchor=east, text=accent] at (-0.12,0) {index rule};
  \node[note, anchor=west, text=accent] at (6.62,0) {$\frac{13}{2}$};
  % fastest: B, C, A
  \draw[hot] (0,-1.25) -- (6.4727,-1.25);
  \foreach \t/\name/\dx in {5/B/0,6.2/C/-0.04,6.4727/A/0.07} {
    \node[circle, draw=warm, fill=white, line width=0.8pt, inner sep=1.5pt] at (\t,-1.25) {};
    \node[text=warm, font=\sffamily\small\bfseries] at (\t+\dx,-0.89) {\name};
  }
  \node[note] at (2.5,-1.57) {income 1};
  \node[note] at (5.6,-1.57) {income 5};
  \node[note, anchor=east, text=warm] at (-0.12,-1.25) {fastest};
  \node[note, anchor=west, text=warm] at (6.62,-1.25) {$\frac{356}{55}$};
\end{tikzpicture}
\caption{Example~\ref{ex:three}. The index rule buys the cheap item $A$ first. The fastest
schedule waits two seconds longer for its first purchase, and it ends $\frac{3}{110}$ seconds
earlier.}\label{fig:three}
\end{figure}

\subsection{A finite list}

\begin{lemma}[no wait helps]\label{lem:earliest}
Let all gains be nonnegative. In a schedule, each purchase is at the same time as in the
earliest schedule of the same order, or later.
\end{lemma}

\begin{proof}
Let the schedule have the times $t_1\le\dots\le t_n$, and let the earliest schedule of its order
have the times $e_1\le\dots\le e_n$. Let $k\ge0$, suppose that $e_j\le t_j$ for all $j\le k$, and
suppose that $e_{k+1}>t_{k+1}$.

Before the time $t_{k+1}$, the schedule has made some of its first $k$ purchases and no others.
The earliest schedule has made each of these purchases at the same time or earlier, and it
makes purchase $k+1$ after $t_{k+1}$. The gains are nonnegative, so at each moment before
$t_{k+1}$ its income is at least the income of the schedule. Thus up to the time $t_{k+1}$ it has
received at least as many cookies. At that time it has paid the first $k$ prices and no others,
and the schedule has paid these prices also. So its bank at the time $t_{k+1}$ holds at least
the bank of the schedule before purchase $k+1$, and that bank holds the price of item $k+1$.
Then the earliest schedule buys item $k+1$ at the time $t_{k+1}$ or earlier. This is not
possible, so $e_{k+1}\le t_{k+1}$. The lemma follows by induction on $k$.
\end{proof}

\begin{theorem}[a finite list]\label{thm:list}
Let $R$ be a finite list of items with prices $p_i\ge0$ and gains $\Delta_i\ge0$. For a bank
$b\ge0$ and an income $c>0$ put
\[
  w_i=\frac{\max(p_i-b,\,0)}{c},\qquad b_i=\max(b-p_i,\,0),
\]
and define
\[
  V(b,c,\emptyset)=0,\qquad
  V(b,c,R)=\min_{i\in R}\Bigl[\,w_i+V\bigl(b_i,\ c+\Delta_i,\ R\setminus\{i\}\bigr)\Bigr].
\]
Then $V(b,c,R)$ is the least time of a schedule that buys all items of $R$.
\end{theorem}

\begin{proof}
By Lemma~\ref{lem:earliest} the least time is the least time of an earliest schedule. We use
induction on the number of items. Let an earliest schedule have the first item $i$. It buys $i$
after $w_i$ seconds. The bank is then $b_i$ and the income is $c+\Delta_i$. From this moment the
schedule is an earliest schedule for the items of $R\setminus\{i\}$, and each earliest schedule
for these items can follow. By the induction hypothesis the least time of the rest is
$V(b_i,c+\Delta_i,R\setminus\{i\})$. The minimum over the first item $i$ is $V(b,c,R)$.
\end{proof}

The state after the purchases of a set $S$ of items depends on $S$ only, and not on the order.
The income is $c+\sum_{i\in S}\Delta_i$. The bank is $\max(b-\sum_{i\in S}p_i,\,0)$, because
$\max(\max(b-p,0)-q,\,0)=\max(b-p-q,\,0)$ for $q\ge0$. So a program needs one value of $V$ for
each subset of $R$. For Example~\ref{ex:three} the recursion gives $\frac{356}{55}$.

The model has a fixed list of items. In the game, a purchase can change other prices and gains,
and a building has many copies. There the recursion is not available and the index rule is a
heuristic. Example~\ref{ex:three} shows a loss of $0.4\%$. We do not know a bound for the loss
in general.

\section{The bank for golden cookies}\label{sec:bank}

With golden cookies, a cookie in the bank has a use. A Lucky golden cookie pays
\begin{equation}\label{eq:lucky}
  g\,\min(0.15\,b,\ 900\,c)+13
\end{equation}
cookies. Here $b$ is the bank, $c$ is the income, and $g\ge1$ is a factor from upgrades. The
payout increases with the bank up to the \emph{cap} $b=6000\,c$ (Figure~\ref{fig:lucky}).

\begin{figure}[ht]
\centering
\begin{tikzpicture}[x=0.00105cm,y=0.0036cm]
  \draw[edge, slate!70, -{Stealth[length=4pt]}] (0,0) -- (9700,0);
  \draw[edge, slate!70, -{Stealth[length=4pt]}] (0,0) -- (0,1120);
  \foreach \x in {3000,6000,9000} {
    \draw[edge, slate!70] (\x,0) -- (\x,-22);
    \node[note] at (\x,-75) {\x};
  }
  \foreach \y in {450,900} {
    \draw[edge, slate!70] (0,\y) -- (-70,\y);
    \node[note, anchor=east] at (-110,\y) {\y};
  }
  \draw[edge, slate!45, dashed] (6000,0) -- (6000,900);
  \draw[edge, slate!45, dashed] (0,900) -- (6000,900);
  \draw[cool] (0,0) -- (6000,900);
  \draw[hot] (6000,900) -- (9300,900);
  \node[port] at (6000,900) {};
  \node[note, text=accent, anchor=south east] at (2950,480) {slope 0.15};
  \node[note, text=warm, anchor=south] at (7650,930) {no gain above the cap};
  \node[note, anchor=north] at (4850,-150) {bank, in seconds of income};
  \node[note, anchor=south, rotate=90] at (-900,500) {payout, in seconds of income};
\end{tikzpicture}
\caption{The Lucky payout \eqref{eq:lucky} for $g=1$, without the constant $13$. Below the cap,
each banked cookie adds $0.15$ cookies to the payout. The cap is at a bank of $6000$ seconds of
income.}\label{fig:lucky}
\end{figure}

\paragraph{Model.} Lucky cookies arrive at the rate $\lambda$: in a short time $dt$ the
expected number of Lucky cookies is $\lambda\,dt$. The player clicks each one at once. The
\emph{expected income} at the bank $b$ and the income $c$ is
\[
  R(b,c)=c+\lambda\bigl(g\,\min(0.15\,b,\ 900\,c)+13\bigr).
\]

\begin{maintheorem}{the bank}\label{thm:bank}
Let an item have a price $p\le b$ and a gain $\Delta\ge0$. The purchase changes the expected
income by
\[
  R(b-p,\,c+\Delta)-R(b,c)=\Delta+\lambda g\Bigl[\min\bigl(0.15\,(b-p),\ 900\,(c+\Delta)\bigr)
  -\min\bigl(0.15\,b,\ 900\,c\bigr)\Bigr].
\]
\begin{enumerate}[leftmargin=1.9em,label=\textup{(\alph*)}]
\item If $b\le6000\,c$, the change is $\Delta-0.15\,\lambda g\,p$. It is positive exactly when
\[
  \frac{p}{\Delta}<\theta,\qquad\theta=\frac{1}{0.15\,g\lambda}.
\]
\item If $b-p\ge6000\,(c+\Delta)$, the change is $(1+900\,\lambda g)\,\Delta$, which is not
negative.
\end{enumerate}
\end{maintheorem}

\begin{proof}
The first formula is the definition of $R$. In case~(a) we have
$0.15\,(b-p)\le0.15\,b\le900\,c\le900\,(c+\Delta)$. So each minimum is its first term, and the
bracket is $-0.15\,p$. In case~(b) we have $0.15\,b\ge0.15\,(b-p)\ge900\,(c+\Delta)\ge900\,c$.
So each minimum is its second term, and the bracket is $900\,\Delta$.
\end{proof}

The number $\theta$ is the time in which a banked cookie earns one cookie. Below the cap, the
theorem says: buy an item if its payback time is less than $\theta$, and keep the cookies if it
is more. Above the cap, a purchase that leaves the bank above the new cap loses nothing. Between
the two cases, the first formula of the theorem gives the change.

The theorem compares two expected incomes at one moment. It does not give the fastest schedule
for a list of purchases (Section~\ref{sec:intro}).

\subsection{The rate of Lucky cookies}

Let $I$ be the time between two golden cookies, in seconds, and let $\pi$ be the share of the
golden cookies that give Lucky. Then $\lambda=\pi/\EE I$.

\begin{lemma}[the time between golden cookies]\label{lem:interval}
After a golden cookie the game waits $t_{\min}$ frames. In each later frame $j=1,\dots,W$ a
golden cookie appears with chance $(j/W)^{5}$, if none has appeared before. Let $J$ be the
number of the frame in which the golden cookie appears. Then
\[
  \EE J=\sum_{k=0}^{W-1}\ \prod_{j=1}^{k}\Bigl(1-\frac{j^{5}}{W^{5}}\Bigr)
  \ \le\ 1+\Gamma\bigl(\tfrac76\bigr)\,\bigl(6\,W^{5}\bigr)^{1/6}.
\]
\end{lemma}

\begin{proof}
The chance that no golden cookie appears in the frames $1,\dots,k$ is
$q_k=\prod_{j\le k}(1-j^{5}/W^{5})$, and $q_W=0$. So $\EE J=\sum_{k\ge0}\PP(J>k)=\sum_{k=0}^{W-1}q_k$.
For the bound we use $1-x\le e^{-x}$ and $\sum_{j\le k}j^{5}\ge\int_0^{k}x^{5}\,dx=k^{6}/6$. They
give $q_k\le f(k)$ with $f(x)=\exp(-x^{6}/(6W^{5}))$. The function $f$ decreases, so
\[
  \sum_{k\ge0}f(k)\le f(0)+\int_0^{\infty}f(x)\,dx=1+\bigl(6\,W^{5}\bigr)^{1/6}\,\Gamma\bigl(\tfrac76\bigr).
\]
The last step is the integral $\int_0^{\infty}e^{-ax^{6}}\,dx=a^{-1/6}\,\Gamma(\frac76)$.
\end{proof}

The game has 30 frames in a second. Without upgrades, $t_{\min}$ is 5 minutes and
$W=t_{\max}-t_{\min}$ is 10 minutes. The upgrades Lucky day and Serendipity each halve the two
numbers. The bound is larger than the exact mean by less than two frames
(Table~\ref{tab:interval}). If the player clicks each golden cookie at once, then $\EE I=(t_{\min}+\EE J)/30$.

\begin{table}[ht]
\centering\small
\begin{tabular}{@{}lrrrrr@{}}
\toprule
Upgrades & $W$ & $\EE J$ & bound & $\EE I$ in seconds & $\theta$ in seconds \\
\midrule
none & 18\,000 & 4396.85 & 4398.13 & 446.56 & 7060 \\
one & 9\,000 & 2467.53 & 2468.80 & 232.25 & 3672 \\
both & 4\,500 & 1384.74 & 1386.01 & 121.16 & 1915 \\
\bottomrule
\end{tabular}
\caption{The time between golden cookies by Lemma~\ref{lem:interval}, and the number $\theta$ of
Theorem~\ref{thm:bank} for $g=1$ and $\pi=0.4217$.}\label{tab:interval}
\end{table}

\begin{fineprint}
The share $\pi=0.4217$ is for golden cookies without wrath and a player with at least ten
buildings. The game makes a list with Frenzy and Lucky. It adds Cookie chain and Cookie storm
together with chance $0.03$, Click frenzy with chance $0.1$, and a building effect with chance
$0.25$. With chance $0.8$ it removes the effect of the last golden cookie from the list. Then it
takes one effect from the list, each with the same chance. The number $0.4217$ is the long-run
share of Lucky in this chain. We leave out effects with a chance below $0.001$ and the effects
of the dragon.
\end{fineprint}

With the two upgrades, $\theta$ is about 32 minutes. Early in a run the payback times are
seconds or minutes, and the rule says to spend. Late in a run they are hours, and the rule says
to keep the bank at the cap.

\section{The order of permanent upgrades}\label{sec:chips}

\paragraph{Model.} There are $n$ upgrades, and the player wants all of them. Upgrade $j$ costs
$p_j>0$ chips and multiplies the speed by $m_j>1$. Chips arrive one after another. At the speed
$S$, a chip takes $q/S$ seconds, with a fixed number $q>0$. At the start the speed is $S_0$ and
the player has no chips. The player buys the upgrades in a fixed order, each at the moment at
which the chips for it are there. So the time of the order $j_1,\dots,j_n$ is
\begin{equation}\label{eq:chiptime}
  \frac{q}{S_0}\sum_{k=1}^{n}\frac{p_{j_k}}{m_{j_1}\cdots m_{j_{k-1}}}.
\end{equation}

\begin{maintheorem}{the ratio rule}\label{thm:ratio}
Put
\[
  \rho_j=\frac{1-1/m_j}{p_j}.
\]
An order has the least time \eqref{eq:chiptime} if and only if $\rho$ does not increase along
it.
\end{maintheorem}

\begin{proof}
Let $i$ and $j$ be adjacent in an order, with $i$ first, and let $\Pi$ be the product of the
factors of the upgrades before them. An exchange of $i$ and $j$ does not change the other terms
of \eqref{eq:chiptime}, because the product after the two upgrades is $\Pi\,m_im_j$ in each
order. The two terms change from $p_i+p_j/m_i$ to $p_j+p_i/m_j$, each times $q/(S_0\Pi)$. The
difference is
\[
  \Bigl(p_j+\frac{p_i}{m_j}\Bigr)-\Bigl(p_i+\frac{p_j}{m_i}\Bigr)
  =p_j\Bigl(1-\frac1{m_i}\Bigr)-p_i\Bigl(1-\frac1{m_j}\Bigr)=p_i\,p_j\,(\rho_i-\rho_j).
\]
So the exchange makes the time less if $\rho_i<\rho_j$, and it does not change the time if
$\rho_i=\rho_j$.

There are finitely many orders, so some order has the least time. An order with $\rho_i<\rho_j$
for an adjacent pair is not such an order. So $\rho$ does not increase along an order with the
least time. Two orders along which $\rho$ does not increase differ by exchanges of adjacent
upgrades with the same $\rho$. So all of them have the same time, which is the least time.
\end{proof}

The proof is the exchange argument of Smith's rule for the order of jobs on one
machine~\cite{Smith}.

The rule can tell the player to save chips. Let $X$ cost 10 chips with the factor $1.05$, and
let $Y$ cost 100 chips with the factor $3$. Then $\rho_X=0.0048$ and $\rho_Y=0.0067$. With $Y$
first the time is $100+10/3=103.3$ units of $q/S_0$, and with $X$ first it is
$10+100/1.05=105.2$ units. So the player keeps the first 10 chips and waits for $Y$.

Some upgrades need other upgrades first. Suppose that a group of upgrades stays together, in a
fixed order $1,\dots,r$. In the exchange argument the group acts as one upgrade with the factor
$M=m_1\cdots m_r$ and the price
\[
  d=\sum_{k=1}^{r}\frac{p_k}{m_1\cdots m_{k-1}},
\]
because two adjacent groups have the times $d+d'/M$ and $d'+d/M'$ in the two orders. The ratio
of the group is $(1-1/M)/d$. The price $d$ is less than the sum of the prices.

\begin{example}[the effort for a chip must be fixed]\label{ex:effort}
Let the chip number $B$ take $B/S$ seconds: the effort increases. In the sum over the chips we
write an integral over $B$. Upgrade $X$ costs 1 chip with the factor $2$, and upgrade $Y$ costs
2 chips with the factor $100$. Then $\rho_X=\frac12$ is more than $\rho_Y=\frac{99}{200}$. With
$S_0=1$, the order with $X$ first has the time
\[
  \int_0^{1}B\,dB+\frac12\int_1^{3}B\,dB=\frac12+2=\frac52,
\]
and the order with $Y$ first has the time
\[
  \int_0^{2}B\,dB+\frac1{100}\int_2^{3}B\,dB=2+\frac1{40}=\frac{81}{40}.
\]
So the order of the ratio rule is the slower one.
\end{example}

In the game the effort is not fixed. By \eqref{eq:prestige}, the chip number $B$ needs about
$3\cdot10^{12}\,B^{2}$ cookies. So the ratio rule is exact in its model and a heuristic in the
game.

\section{The time of a reset}\label{sec:reset}

\subsection{Speed}

Permanent bonuses multiply the income of a run. A run has the \emph{speed} $S$ if each source of
income is $S$ times the source in a run with the speed $1$, with the same items.

\begin{lemma}[time and speed]\label{lem:rescale}
Suppose that the decisions of the player depend on the bank and on the items, and not on the
clock. If the run with the speed $1$ bakes $F(x)$ cookies in its first $x$ seconds, then the run
with the speed $S$ bakes $F(St)$ cookies in its first $t$ seconds.
\end{lemma}

\begin{proof}
Put $s=St$. Between two purchases, the bank of the run with the speed $S$ increases by
$S\,r\,dt=r\,ds$ in the time $dt$, where $r$ is the income at the speed $1$ with the same items.
The prices are the same, and the decisions see the same bank and the same items. So in the
variable $s$ the run is the run with the speed $1$.
\end{proof}

In the game the lemma is exact for the buildings and the clicks. Golden cookies follow the
clock, so for them it is an approximation.

\begin{proposition}[a pause for a permanent gain]\label{prop:pause}
Suppose that a task needs $H$ more seconds at the present speed. An action stops all income for
$D\ge0$ seconds and then multiplies the speed by $M>0$. With the action the task needs $D+H/M$
seconds. So the action shortens the task exactly when
\[
  D<H\Bigl(1-\frac1M\Bigr).
\]
\end{proposition}

\begin{proof}
By Lemma~\ref{lem:rescale}, a speed that is $M$ times larger divides the remaining time by $M$.
So the action gives the time $D+H/M$, and $D+H/M<H$ is the stated inequality.
\end{proof}

For $M=1+x$ the action can cost up to $Hx/(1+x)$ seconds. This is less than $Hx$. In the game,
an achievement is such a gain. In one game with eight kitten upgrades, one more achievement gave
$x=0.0214$. For a task of 77 hours, an achievement was then worth a pause of up to
$77\cdot0.0214/1.0214=1.61$ hours.

\subsection{A chain of resets}

\paragraph{Model.} The prestige is the real number $P$ of \eqref{eq:prestige}. A run that
starts with the prestige $P$ has the speed $\kappa P$, with a fixed number $\kappa>0$. A run
with the speed $1$ bakes $F(x)=\alpha\,x^{\beta}$ cookies in its first $x$ seconds, with
$\alpha>0$ and $\beta>1$. A reset takes no time. The earlier runs have baked $A_0>0$ cookies, and the goal is
$A^{\ast}>A_0$. A \emph{schedule} is a chain
$A_0<A_1<\dots<A_n=A^{\ast}$: run number $k$ starts at $A_{k-1}$ and the player resets at $A_k$.

A run from $A$ to $vA$, with $v>1$, bakes $(v-1)A$ cookies. By Lemma~\ref{lem:rescale} it takes
the time
\begin{equation}\label{eq:runtime}
  T(A,v)=\frac{\bigl((v-1)A/\alpha\bigr)^{1/\beta}}{\kappa\,(A/10^{12})^{1/3}}
  =K\,A^{\gamma/3}\,(v-1)^{1/\beta},\qquad
  \gamma=\frac3\beta-1,\quad K=\frac{10^{4}}{\kappa\,\alpha^{1/\beta}}.
\end{equation}
For $\beta<3$ the exponent $\gamma$ is positive, and later runs take longer for the same factor
$v$. For $\beta>3$ they take less time.

\begin{maintheorem}{the time of a reset}\label{thm:reset}
Put $\Phi(A)=(3/\gamma)\,A^{\gamma/3}$ if $\gamma\ne0$ and $\Phi(A)=\ln A$ if $\gamma=0$. For
$v>1$ put
\[
  \rho(v)=\frac{(v-1)^{1/\beta}}{\Phi(v)-\Phi(1)}.
\]
\begin{enumerate}[leftmargin=1.9em,label=\textup{(\alph*)}]
\item The function $\rho$ decreases on an interval $(1,v^{\star}]$ and increases on
$[v^{\star},\infty)$. The number $v^{\star}$ is the only solution $v>1$ of
\begin{equation}\label{eq:vstar}
  \frac\beta3\Bigl(1-\frac1v\Bigr)=\frac{1-v^{-\gamma/3}}{\gamma}\quad(\gamma\ne0),
  \qquad\qquad 3\,(v-1)=v\ln v\quad(\gamma=0).
\end{equation}
\item Each schedule takes at least the time $K\,\rho(v^{\star})\,\bigl(\Phi(A^{\ast})-\Phi(A_0)\bigr)$.
\item A schedule takes exactly this time if and only if $A_k=v^{\star}A_{k-1}$ for each $k$.
\end{enumerate}
\end{maintheorem}

\begin{proof}
\step{The sum.} The function $\Phi$ increases, and $\Phi(vA)-\Phi(A)=A^{\gamma/3}\,(\Phi(v)-\Phi(1))$
for all $A>0$ and $v>1$. So \eqref{eq:runtime} is
\[
  T(A,v)=K\,\rho(v)\,\bigl(\Phi(vA)-\Phi(A)\bigr).
\]
Let a schedule have the factors $v_k=A_k/A_{k-1}$. Its time is
\begin{align*}
  \sum_{k=1}^{n}K\,\rho(v_k)\,\bigl(\Phi(A_k)-\Phi(A_{k-1})\bigr)
  &\ \ge\ K\,\min_{v>1}\rho(v)\ \sum_{k=1}^{n}\bigl(\Phi(A_k)-\Phi(A_{k-1})\bigr)\\
  &\ =\ K\,\min_{v>1}\rho(v)\,\bigl(\Phi(A^{\ast})-\Phi(A_0)\bigr).
\end{align*}
The two sides are equal if and only if each $v_k$ is a minimum point of $\rho$. So (b) and (c)
follow from (a).

\step{One minimum.} Put $u=\ln v$, and $\psi(x)=x/(1-e^{-x})$ for $x\ne0$, $\psi(0)=1$. The
function $\psi$ is positive. A calculation gives
\[
  \frac{d}{du}\ln\rho=\frac1\beta\cdot\frac{e^{u}}{e^{u}-1}
  -\frac\gamma3\cdot\frac{e^{\gamma u/3}}{e^{\gamma u/3}-1}
  =\frac1u\Bigl(\frac{\psi(u)}{\beta}-\psi\bigl(\tfrac{\gamma u}{3}\bigr)\Bigr),
\]
where the second term is $1/u$ if $\gamma=0$. So the derivative has the sign of
\[
  h(u)=\ln\psi(u)-\ln\psi\bigl(\tfrac{\gamma u}{3}\bigr)-\ln\beta .
\]
Put $L(x)=x\,\psi'(x)/\psi(x)=1-x/(e^{x}-1)$, with $L(0)=0$. The function $x/(e^{x}-1)$ has a
derivative with the numerator $e^{x}(1-x)-1$. This numerator has its maximum $0$ at $x=0$, so
it is negative for $x\ne0$. So $L$ increases strictly on the real line. We have $\gamma/3<1$, and thus
\[
  u\,h'(u)=L(u)-L\bigl(\tfrac{\gamma u}{3}\bigr)>0\qquad(u>0).
\]
So $h$ increases strictly. For $u\to0$ we have $h(u)\to-\ln\beta<0$. For $u\to\infty$ we have
$\psi(u)/u\to1$. If $\gamma\le0$, then $\psi(\gamma u/3)\le1$ and $h(u)\to\infty$. If $\gamma>0$,
then $\psi(\gamma u/3)/u\to\gamma/3$ and $h(u)\to\ln(3/(\beta\gamma))=\ln(3/(3-\beta))>0$. Hence
$h$ has exactly one zero $u^{\star}$. It is negative before $u^{\star}$ and positive after it,
and (a) holds with $v^{\star}=e^{u^{\star}}$. The zero satisfies
$\psi(u)=\beta\,\psi(\gamma u/3)$, which is \eqref{eq:vstar}.
\end{proof}

So the player resets when the cookies of all runs are $v^{\star}$ times their number at the
start of the run. The prestige is then $(v^{\star})^{1/3}$ times its value at the start
(Table~\ref{tab:vstar}). For $\beta=\frac32$ the equation \eqref{eq:vstar} is
$x^{3}-2x^{2}+1=0$ for $x=v^{1/3}$, and the factor of the prestige is the golden ratio
$(1+\sqrt5)/2$.

If $A^{\ast}/A_0$ is not a power of $v^{\star}$, then no schedule has the time of (b), and the
least time is larger.

\begin{table}[ht]
\centering\small
\begin{tabular}{@{}rrrrr@{}}
\toprule
$\beta$ & $\gamma$ & $v^{\star}$ & $(v^{\star})^{1/3}$ & root of \eqref{eq:mvt} \\
\midrule
1.5 & 1 & 4.24 & 1.62 & 2.40 \\
2 & 0.5 & 7.87 & 1.99 & 4.92 \\
2.5 & 0.2 & 12.03 & 2.29 & 9.31 \\
3 & 0 & 16.80 & 2.56 & 16.80 \\
4 & $-0.25$ & 28.39 & 3.05 & 50.44 \\
6 & $-0.5$ & 60.91 & 3.93 & 397.38 \\
\bottomrule
\end{tabular}
\caption{The factor $v^{\star}$ of Theorem~\ref{thm:reset} for six growth exponents, the factor
of the prestige, and the factor of the rule without the correction for later
runs.}\label{tab:vstar}
\end{table}

The exponent $\beta$ of the present run has an estimate. If the run has baked $C$ cookies in
$t$ seconds and the income is $c$ now, then $F(x)=\alpha x^{\beta}$ gives $\beta=c\,t/C$.

For $\gamma=0$ the theorem says: reset at the moment at which $\ln P$ increases at its mean rate
over the run. This is the marginal value theorem of Charnov~\cite{Charnov}. For
$\beta\ne3$ the same rule ignores that later runs take more or less time. It then has the
equation
\begin{equation}\label{eq:mvt}
  \beta\,(v-1)=v\ln v
\end{equation}
in the place of \eqref{eq:vstar} (Figure~\ref{fig:vstar}).

\begin{figure}[ht]
\centering
\begin{tikzpicture}[x=1.55cm,y=1.35cm]
  % frame and grid; the vertical axis is log10 of the factor
  \foreach \y/\lab in {0/1,1/10,2/100,3/1000} {
    \draw[edge, slate!25] (1,\y) -- (8.1,\y);
    \node[note, anchor=east] at (0.93,\y) {\lab};
  }
  \foreach \x in {2,...,8} {
    \draw[edge, slate!70] (\x,0) -- (\x,-0.06);
    \node[note] at (\x,-0.24) {\x};
  }
  \draw[edge, slate!70] (1,0) -- (8.1,0);
  \draw[edge, slate!70] (1,0) -- (1,3.55);
  \draw[edge, slate!45, dashed] (3,0) -- (3,1.2253);
  \draw[hot, dashed] plot coordinates {(1.2,0.1635) (1.4,0.3107) (1.6,0.4461) (1.8,0.5726) (2.0,0.6921) (2.2,0.8061) (2.4,0.9158) (2.6,1.0218) (2.8,1.1248) (3.0,1.2253) (3.2,1.3238) (3.4,1.4205) (3.6,1.5158) (3.8,1.6098) (4.0,1.7027) (4.2,1.7948) (4.4,1.8861) (4.6,1.9767) (4.8,2.0667) (5.0,2.1563) (5.2,2.2455) (5.4,2.3343) (5.6,2.4229) (5.8,2.5111) (6.0,2.5992) (6.2,2.6871) (6.4,2.7748) (6.6,2.8624) (6.8,2.9499) (7.0,3.0373) (7.2,3.1246) (7.4,3.2118) (7.6,3.2990) (7.8,3.3861) (8.0,3.4732)};
  \draw[cool] plot coordinates {(1.2,0.3522) (1.4,0.5515) (1.6,0.6926) (1.8,0.8035) (2.0,0.8961) (2.2,0.9762) (2.4,1.0474) (2.6,1.1117) (2.8,1.1707) (3.0,1.2253) (3.2,1.2763) (3.4,1.3242) (3.6,1.3694) (3.8,1.4123) (4.0,1.4532) (4.2,1.4922) (4.4,1.5296) (4.6,1.5655) (4.8,1.6001) (5.0,1.6334) (5.2,1.6656) (5.4,1.6968) (5.6,1.7270) (5.8,1.7563) (6.0,1.7847) (6.2,1.8124) (6.4,1.8393) (6.6,1.8655) (6.8,1.8911) (7.0,1.9160) (7.2,1.9404) (7.4,1.9642) (7.6,1.9875) (7.8,2.0102) (8.0,2.0325)};
  \node[port] at (3,1.2253) {};
  \node[note, text=accent, anchor=west] at (6.1,1.52) {$v^{\star}$ of \eqref{eq:vstar}};
  \node[note, text=warm, anchor=east] at (5.85,2.76) {root of \eqref{eq:mvt}};
  \node[note, anchor=west] at (3.08,0.42) {the two agree at $\beta=3$};
  \node[note, anchor=north] at (4.55,-0.42) {growth exponent $\beta$};
  \node[note, anchor=south, rotate=90] at (0.28,1.75) {factor of all cookies in one run};
\end{tikzpicture}
\caption{The factor $v^{\star}$ and the root of \eqref{eq:mvt}. The vertical scale is
logarithmic.}\label{fig:vstar}
\end{figure}

\subsection{Early in the game}

\begin{proposition}[no gain in speed, no reset]\label{prop:early}
Suppose that all runs have the same speed $S$, that $F(x)=\alpha\,x^{\beta}$ with $\beta>1$, and
that a reset takes $\tau\ge0$ seconds. Then one run is the fastest way to bake a given number
of cookies.
\end{proposition}

\begin{proof}
Let $n\ge2$ runs bake $D_1,\dots,D_n>0$ cookies, with the sum $D$. Their time is
\[
  \sum_{k=1}^{n}\frac{(D_k/\alpha)^{1/\beta}}{S}+(n-1)\,\tau .
\]
For $0<a<1$ and positive numbers we have $(x+y)^{a}<x^{a}+y^{a}$. With $a=1/\beta$ this gives
$D^{1/\beta}<\sum_kD_k^{1/\beta}$. So one run with $D$ cookies takes less time.
\end{proof}

Early in the game a reset changes the speed by a small amount, and the game is near this case.
Theorem~\ref{thm:reset} is for the late game, where the bonus of the prestige is much larger
than the income without it.

With a reset time $\tau>0$ or a different speed, we know no formula. A program can take a finite
set of levels $a_0<a_1<\dots<a_n=A^{\ast}$ and the time $t(i,j)$ of a run from $a_i$ to $a_j$.
The least time $V(i)$ from $a_i$ with resets at these levels satisfies
\[
  V(n)=0,\qquad V(i)=\min_{j>i}\bigl[\,t(i,j)+\tau\,[j<n]+V(j)\,\bigr],
\]
where $[j<n]$ is $1$ if $j<n$ and $0$ if $j=n$. The reason is that a schedule is a path from
$a_i$ to $a_n$, and its first run ends at some $a_j$. The choice $j=n$ is the run that reaches
the goal without a reset. A program must include it.

\begin{thebibliography}{9}\small

\bibitem{Charnov} E.~L.~Charnov, Optimal foraging, the marginal value theorem,
\emph{Theoret.\ Population Biol.}\ 9 (1976), 129--136.

\bibitem{CM} Cookie Monster Team, Cookie Monster, an add-on for Cookie Clicker,
\url{https://github.com/CookieMonsterTeam/CookieMonster}.

\bibitem{Repo} KNT and Fable~5.2, a guide to Cookie Clicker, the mod Cookie Monster+ and the
Lean files of this paper, \url{https://jehlp.net/cookie-clicker/} and
\url{https://github.com/JWKNT/cookie-clicker}.

\bibitem{Smith} W.~E.~Smith, Various optimizers for single-stage production,
\emph{Naval Res.\ Logist.\ Quart.}\ 3 (1956), 59--66.

\bibitem{Game} J.~Thiennot (Orteil), Cookie Clicker, version 2.053, DashNet,
\url{https://orteil.dashnet.org/cookieclicker/}.

\end{thebibliography}

\end{document}
